\usepackage{shellesc}
\usepackage{polyglossia}
	\setmainlanguage{french}
\usepackage{tikz} 
	\usetikzlibrary{shapes}
	\usetikzlibrary{arrows, arrows.meta}
	\usetikzlibrary{calc}
	\usetikzlibrary{babel}
	\usetikzlibrary{patterns}
	\usetikzlibrary{matrix}
	\usetikzlibrary{backgrounds}
	\usetikzlibrary{external}
	\usetikzlibrary{fit} %positioning, 
	\usetikzlibrary{positioning}
	\usetikzlibrary{intersections}
	\usetikzlibrary{tikzmark}
%
\usepackage{pgfplotstable}
\usepackage[
	separate-uncertainty=true,
	multi-part-units=single,
	]{siunitx}
\usepackage{ulem}
%
\usepackage{animate}
\usepackage{pgfplots}
\pgfplotsset{compat=1.8}
\usepackage{xifthen} %for conditional test
%
%\usepackage[table,xcdraw]{xcolor}
\usepackage{pdfpages}
\usepackage[absolute,overlay]{textpos}
\usepackage{graphicx}
\usepackage{capt-of}
\usepackage{hyperref} % for online version
\usepackage[
	style = authoryear,
	maxcitenames = 1,
	maxbibnames = 99,
	backref = true,
	hyperref = true,
	uniquelist = false, %to avoid more than 2 authors in text
	uniquename = false,
	backend = biber,
	eprint = false,
	url = false,
	giveninits = true,
]{biblatex}
\usepackage{cleveref} % should be loaded last

\usetheme{Frankfurt}
%\useinnertheme[shadow=false]{rounded}

\definecolor{Prune}{RGB}{99,0,60}
\definecolor{LimedSpruce}{RGB}{49,62,72}
\definecolor{Shiraz}{RGB}{198,11,70}
\definecolor{Orient}{RGB}{0,78,125}
\definecolor{Teal}{RGB}{0,128,122}
\definecolor{Pear}{RGB}{213,223,61}
\definecolor{Lavender}{RGB}{185,177,215}
\definecolor{Sushi}{RGB}{140,198,62}
\definecolor{Pistachio}{RGB}{157,194,9}
\definecolor{Jewel}{RGB}{18,107,64}

\colorlet{DarkPrune}{Prune!50!black}
\colorlet{DarkOrient}{Orient!50!black}
\colorlet{DarkShiraz}{Shiraz!50!black}

\setbeamercolor{title}{bg=Prune, fg=white}
\setbeamercolor{frametitle}{bg=Prune, fg=white}
\setbeamercolor{section in head/foot}{fg=white, bg=DarkPrune}
\setbeamertemplate{itemize items}[circle]
\setbeamercolor{itemize item}{fg=Prune}
\setbeamercolor{itemize subitem}{fg=Orient}
\setbeamercolor{block title}{bg=Prune}
\setbeamercovered{transparent}
%\colorlet{bullet}{Prune}

%\setbeamertemplate{titlepage}[default][shadow=false]
\setbeamertemplate{frametitle}[default][shadow=false]
\setbeamertemplate{blocks}[rounded][shadow=false]
\setbeamertemplate{headline}[shadow=false]
\setbeamertemplate{subsection in head}[shadow=false]
\setbeamertemplate{section in head}[shadow=false]
\setbeamertemplate{beamercolorbox}[shadow=false]

\setbeamertemplate{footline}[frame number] % start slides numbering

\setbeamertemplate{section in toc}[sections numbered]
\setbeamercolor{section in toc}{fg=Prune}
\setbeamertemplate{subsection in toc}[subsections numbered]

\addbibresource{../bibliography/introduction.bib}
\addbibresource{../bibliography/chapter1.bib}
\addbibresource{../bibliography/chapter2.bib}
\addbibresource{../bibliography/chapter3.bib}
\addbibresource{../bibliography/chapter4.bib}
\addbibresource{../bibliography/chapter5.bib}

%\hypersetup{
%	% draft=true, % for printable version
%	colorlinks=true,
%	linkcolor=DarkOrient,
%	filecolor=Sushi,
%	citecolor=DarkShiraz,      
%	urlcolor=Prune,
%}

\AtBeginSection[]
{
	\begin{frame}<beamer>
	\frametitle{\insertsection}
	\tableofcontents[currentsection,sectionstyle=hide/hide,subsectionstyle=show/show/hide]
\end{frame}
}

% for figure opacity
\newcommand<>{\uncovergraphics}[2][{}]{
	% Taken from: <https://tex.stackexchange.com/a/354033/95423>
	\begin{tikzpicture}
		\node[anchor=south west,inner sep=0] (B) at (4,0)
		{\includegraphics[#1]{#2}};
		\alt#3{}{%
			\fill [draw=none, fill=white, fill opacity=0.9] (B.north west) -- (B.north east) -- (B.south east) -- (B.south west) -- (B.north west) -- cycle;
		}
	\end{tikzpicture}
}

\renewcommand<>{\sout}[1]{%
	\alt#2{\beameroriginal{\sout}{#1}}{#1}
}

%% math
\newcommand{\vect}[1]{
	\boldsymbol{\MakeLowercase{#1}}
}
%
\newcommand{\matr}[1]{
	\boldsymbol{\MakeUppercase{#1}}
}
%
\newcommand{\matrCal}[1]{
	\boldsymbol{\mathcal{\MakeUppercase{#1}}}
}
%
\newcommand{\jacobian}[3]{ %
	\ifthenelse{\isempty{#1}}
	{%
		\boldsymbol{\mathcal{J}_{#2}^{#3}} %
	}
	{%
		\boldsymbol{\mathcal{J}_{#2}}(#1)^{#3} %
	}	
}
%
\newcommand{\jacobianHat}[3]{ %
	\ifthenelse{\isempty{#1}}
	{%
		\boldsymbol{\hat{\mathcal{J}}_{#2}^{#3}} %
	}
	{%
		\boldsymbol{\hat{\mathcal{J}}_{#2}}(#1)^{#3} %
	}	
}
%
\newcommand{\jacobianDot}[3]{ %
	\ifthenelse{\isempty{#1}}
	{%
		\boldsymbol{\dot{\mathcal{J}}_{#2}^{#3}} %
	}
	{%
		\boldsymbol{\dot{\mathcal{J}}_{#2}}(#1)^{#3} %
	}	
}
%
\newcommand{\Laplace}{\mathcal{L}}
%
\DeclareMathOperator{\sign}{sgn}
%
\renewcommand{\Re}{\rm I\!R}
%
\newcommand{\Real}{\mathfrak{Re}}
%
%\newcommand{\Im}{\mathfrak{Im}}
%%
\newcommand{\mean}[1]{\overline{#1}}
%
\newcommand{\hadamard}{\odot}
%
\newcommand{\norm}[1]{\left\lVert#1\right\rVert}
%
\newcommand{\abs}[1]{\left\lvert#1\right\rvert}

%%tikz transparency
% Keys to support piece-wise uncovering of elements in TikZ pictures:
% \node[visible on=<2->](foo){Foo}
% \node[visible on=<{2,4}>](bar){Bar}   % put braces around comma expressions
%
% Internally works by setting opacity=0 when invisible, which has the 
% adavantage (compared to \node<2->(foo){Foo} that the node is always there, hence
% always consumes space plus that coordinate (foo) is always available.
%
% The actual command that implements the invisibility can be overriden
% by altering the style invisible. For instance \tikzsset{invisible/.style={opacity=0.2}}
% would dim the "invisible" parts. Alternatively, the color might be set to white, if the
% output driver does not support transparencies (e.g., PS) 
%
\tikzset{
	lowvisible/.style={opacity=0.2},
	visible on/.style={alt={#1{}{lowvisible}}},
	alt/.code args={<#1>#2#3}{%
		\alt<#1>{\pgfkeysalso{#2}}{\pgfkeysalso{#3}} % \pgfkeysalso doesn't change the path
	},
}

\tikzset{
	invisible/.style={opacity=0},
	only on/.style={alt={#1{}{invisible}}},
	alt/.code args={<#1>#2#3}{%
		\alt<#1>{\pgfkeysalso{#2}}{\pgfkeysalso{#3}} % \pgfkeysalso doesn't change the path
	},
}

\graphicspath{{../pdf_figures/}}

\tikzexternalize[optimize command away=\includepdf] % option addedto avoid conflict with pdfpages package
\tikzsetexternalprefix{pdf_figures/}

\tikzexternaldisable